\section{引言}
\label{sec:introduction}
\subsection{EAR：让研究持续推进}
\label{sec:position}
科研中的许多判断，发生在计划改变的时候。刚找到的一篇近邻论文，可能让一个高分选题失去继续做的理由；一个反例，会迫使证明换一条路线；评阅指出的缺口，又会把已经成形的稿件送回推导或实验。研究由这些具体的转折向前推进。自动化系统也需要跟得上：保存已经得到的材料，理解新证据改变了什么，然后继续工作。

EAR（Effective Auto Research）是一套开源科研自动化系统，面向问题发现、自主数学研究和评审回应。我们希望研究者给出方向和约束之后，检索、尝试与修订能够持续进行；再次介入时，有新的产物可以检查，也有清楚的理由说明哪些路线值得继续。

EAR 提供三个可独立使用的工作流。问题发现从方向和资源条件出发，比较文献与候选，形成研究提案；数学研究推进具体目标，交付论证、实现与稿件；评审回应依据论文、审稿意见和已有证据，组织回答并保留依据。研究者可以通过这些材料接续不同阶段的工作。\src{1}

持续推进还需要积累研究做法。在数学研究模块中，EAR 将两者放进一对循环：内循环完成当前研究，外循环用完整研究结果检验执行策略。选题条件、证明组织和修订指令因此有机会随研究经验改进。\textbf{研究结果与研究做法，构成 EAR 关注的两个对象。}系统既要完成当前任务，也要从成功、失败和评阅中获得下一次执行的依据。

\Sub{把有限的人时留给关键判断}
科研中的主动投入分散在阅读、编排、检查和返工中。每次中断之后重新恢复上下文，也要付出额外代价。\src{6,7} EAR 因而把研究者的主动人时置于设计目标的中心：
\begin{equation}
\min H\qquad \text{subject to}\qquad Q\ge Q_{\min},\quad C\le B.
\label{eq:human-objective}
\end{equation}
$H$ 表示主动人时，$Q_{\min}$ 为要求达到的研究质量，$C$ 为模型、工具与算力费用，$B$ 为机器预算。日历周期 $T$ 单独记录，并可作为截止时间约束；多人协作时，$H$ 为所有参与者实际投入的总和。

这个目标要求核算完整的工作负担。检索阅读、任务编排、技术检查、模型错误修复、失败路线返工和最终审阅都计入人时，新增审核工作也必须计入。无人值守的等待属于日历周期，期间没有主动工作的部分不计入 $H$。本文报告过程记录、模型评阅和调用数量，主动人时的实际节省仍待直接测量。

\subsection{持续研究的难点：把反馈变成下一步}
\label{sec:related}
现有工作已经把自动化带进科研的不同环节。ResearchAgent 从文献出发，借助评阅反馈迭代研究点子；The AI Scientist 将点子、实验、结果分析和写作串成科研流程；Paper2Rebuttal 围绕审稿疑虑整理证据与回应计划。GEPA 则把执行轨迹和文本反馈用于提示演化。这些探索为持续研究提供了重要起点。\src{2,3,4,5}

实际使用时，研究者首先面对的是材料的接续。一份提案进入推导之后，它的假设、近邻文献和验证计划仍然有用；一段证明失败之后，反例和失败位置也应当留下。若这些判断散落在不同交互中，接手工作就需要重新阅读和组织。保存怎样的研究状态、如何交接、何时需要人介入，都会影响长期使用的负担。

反馈的接续更进一步。一个低分并没有直接给出下一步：缺口可能在证明，也可能在贡献或表达。系统需要把评阅意见送到相应的研究动作中，让推导得到修复、证据得到补充，或让稿件的主张与已有结果重新对齐。

最后，经验需要进入下一次研究。一条看起来合理的新指令，仍要经历选题、技术尝试、写作和评价，才能知道它是否改善了执行。研究做法的更新，需要由完整任务的结果来决定。表\ref{tab:position} 概括了相关工作各自关注的对象。

\begin{table}[H]
\centering\small
\begin{tabularx}{\textwidth}{P{35mm}Y Y}
\toprule
\rowcolor{warm}
\textbf{工作} & \textbf{研究或优化对象} & \textbf{反馈的作用}\\
\midrule
ResearchAgent\src{2} & 基于文献的问题、方法与实验构想 & 评阅意见用于迭代研究点子\\
The AI Scientist\src{3} & 机器学习点子、实验、结果与论文 & 实验结果支持迭代，自动评阅用于稿件评价与知识档案\\
GEPA\src{4} & 语言模型系统中的提示 & 执行轨迹与反思用于提示变异和组合\\
Paper2Rebuttal\src{5} & 评审疑虑、证据与回应计划 & 先形成可检查计划，再组织回应\\
EAR\src{1,10} & 研究任务及其执行策略 & 当前研究接受反馈，数学外循环用完整研究检验执行策略\\
\bottomrule
\end{tabularx}
\caption{相关工作的研究对象与反馈机制。表中比较设计定位，系统间的同任务性能尚未测试；The AI Scientist 对应 2024 年原始系统，不包含后续 AI Scientist-v2。}
\label{tab:position}
\end{table}

EAR 把这三处接续放在同一设计目标下：研究材料支持继续执行，反馈落实为具体修改，完整研究结果决定后续策略。对于需要反复筛选方向、推进数学与算法任务，或处理评审疑虑的研究者，这种组织方式提供了持续工作和检查关键取舍的基础。

\subsection{EAR 的核心理念：反馈推进研究，结果改进策略}
\label{sec:contributions}
在数学模块中，EAR 把一次完整研究作为策略演化的评价单元。内循环围绕当前目标检索、推导、验证与修订；外循环让候选策略各自执行研究，再依据稿件、评价和失败原因决定保留或修改。一次研究中暴露的问题，可以进入下一次的选题要求、证明组织和修订指令。\textbf{每次研究，都应当为下一次研究留下更好的起点。}

这一理念也体现在另外两个工作流中。问题发现把新文献用于重新判断候选价值：有依据的方案继续精修，发生碰撞或不可行的路线保留退出理由。评审回应先寻找疑虑对应的证据，再补齐缺口并组织回答。反馈因此能够改变投入方向，也能落到具体的技术工作上。

人的判断贯穿其中。研究者确定方向与质量标准，检查关键技术结论，并对最终稿件和回应负责。EAR 承担可持续执行的检索、尝试和修订，同时保留候选取舍、版本变化与策略修改的依据。这种分工希望让研究者每次返回时，都能从已经推进的工作中继续判断。

\begin{insight}
\textbf{从设计主张到可观察的结果}\quad 在 G0--G3 的同一组 28 个完成稿件中，三轮修订后的当前版本模型评阅通过数由 3 增至 12。第\ref{sec:inner-results} 节给出逐轮轨迹与最佳版本结果；这里的通过采用历史模型评阅规则。\src{10}
\end{insight}

下文先展开系统架构与三个工作流，再用稿件修订、策略比较和回应评分预测检查这些机制的表现。
